\documentclass{../styles/cv}

\begin{document}

\populatedtitle

\section{Education}

\subsectionpositiondate
    {York University, Toronto, ON}
    {Master of Science in Geography}
    {Aug 2017 - Feb 2020}
\resumesublistbegin
    \item Thesis: Surface Hydrology Characteristics of Inland Patchy Wetlands in Southern and Western Iceland
    \item Committee: Drs. Kathy L. Young and Taly D. Drezner
    \item Coursework: Dynamics of Snow and Ice \titleseperator Hydrometeorology \titleseperator Ecological Climatology \titleseperator Environmental Policy
\resumesublistend

\subsectionpositiondate
    {DePaul University, Chicago, IL}
    {Bachelor of Arts in Sociology, Minor in Geography: Nature-Society Studies}
    {Sep 2011 - Jun 2015}
\resumesublistbegin
    \item Coursework: Conservation \titleseperator Environmental Justice \titleseperator Soil Science \titleseperator Introductions to GIS and Remote Sensing
\resumesublistend

\section{Fieldwork}

\subsectionpositiondate
    {Inspiring Girls* Expeditions (IGE)}
    {Onsite Coordinator}
    {April 2025 - Present}
\resumesublistbegin
    \item Location: Fairbanks, AK
    \item Attend planning meetings including content: communication styles, education/program goals, leadership style, goals for self growth
    \item Coordinate gear fairy volunteers -- paperwork and lead into the field
\resumesublistend

\subsectionpositiondate
    {Juneau Icefield Research Program (JIRP)}
    {Logistics Support Staff}
    {May 2022 - Present}
\resumesublistbegin
    \item Location: Juneau, AK
    \item Log receipts and track expenses, download forms from Wufoo for processing
    \item Communicate daily with field contact through satellite texting or satellite phone
    \item Order and inventory supplies needed for remote field teams
    \item Collaborate with directors of academics/research and operations to coordinate information across project groups
    \item Coordinate passenger arrival and departure schedules into Juneau and onto the Juneau Icefield
    \item Post program activities and field material on JIRP social media platforms
    \item Schedule helicopter flights and supply runs, adapt schedule to changing weather conditions, and prioritize cargo on flights
    \item Maintain facilities at the Eagle Valley Center and direct visitors
    \item Updated local store accounts with the Foundation for Glacier and Environmental Research to include authorized buyers for 2022
    \item Completed safety training including helicopter safety, Wilderness First Aid, and local hazards in Juneau
    \item Determined program priorities and tracked progress on task-based action items
    \item Adjusted tasks according to priority level, timelines, and time constraints
    \item Coordinated with Directors to procure supplies for field and research teams
    \item Arranged incoming/outgoing travel and administered reimbursements
    \item Supported concurrent Upward Bound program and identified graduate students' teaching strengths
    \item Attended Celebration in support and education of Tlingit, Haida, and Tsimshian cultures
    \item Facilitated wellness and reflection meetings, tracked wellness goals, and led check ins and debriefs
    \item Served as point of contact, responded to program problems/inquiries, and organized pre-field requirements
    \item Implemented daily and weekly schedules to manage tasks and prioritize participant needs
    \item Logged receipts, tracked expenses, downloaded and processed receipts and survey data
\resumesublistend

\section{Teaching and Research}

\subsectionpositiondate
    {York University -- MSc in Arctic Hydrology}
    {Teaching Assistant}
    {Sep 2017 - Feb 2020}
\resumesublistbegin
    \item Location: Toronto, ON
    \item Built relationships with Icelandic professors and landowners to exchange research data
    \item Presented results at national conferences, contributed to publications, and shared on social media
    \item Mentored undergraduate and graduate students (2) on writing, methods, data collection, and arctic safety
    \item Wrote abstracts, grant reports, grant proposals (received \$8,000 in awards) and budgeted for field costs
    \item Planned logistics for arctic fieldwork (sampling locations, shelter and supplies, RV rental, schedule)
    \item Analyzed infiltration and surface soil moisture characteristics in Icelandic drained, patchy wetlands
\resumesublistend

\section{Research Experience}

\subsectionpositiondate
    {York University}
    {MSc Research in Arctic Hydrology}
    {Sep 2017 - Feb 2020}
\resumesublistbegin
    \item Location: Toronto, ON
    \item Supervisor: Dr. Kathy L. Young
    \item Wrote, edited, and presented findings at conferences, thesis defense, and shared results on social media
    \item Assessed and analyzed surface hydrological characteristics of Icelandic drained, patchy wetlands with collected infiltration, near-surface soil moisture, surface albedo, and soils data across 3 wetland soil types
    \item Supervised and mentored an undergraduate student on collecting soil samples in-situ and analyzing soil characteristics in lab (description and characteristics, bulk density, soil organic carbon, soil texture, pH)
    \item Mentored graduate students (2) throughout their program on methods, data collection, arctic safety and travel, grant applications, data analysis, and the writing process
    \item Collected vegetation, meteorological, and landscape data (slope, hummocks) for analysis using HOBO data loggers, Garmin GPS, LANDSAT, ArcticDEM, and a pyranometer and voltmeter
    \item Planned logistics for arctic fieldwork and facilitated communication with professors at the Agricultural University of Iceland to investigate field sites and adapt to sampling scheme as necessary
    \item Coordinated with York University technicians to prepare tools and other collection materials as necessary, and worked on troubleshooting instrument problems in the field
    \item Liaised with wetland landowners and ecological science experts in Iceland to inform project objectives and gain a better understanding of wetland soils
    \item Investigated environmental impacts of ditching on Icelandic wetland habitats, hydrological and ecological processes, and other potential threats to surviving inland patchy wetlands
    \item Researched northern hydrology, peat and volcanic wetland soils, soil moisture, infiltration in frozen/non-frozen soils, and northern climatology as it related to Icelandic patchy wetlands
    \item Wrote and edited grant proposals for research project, budgeted for fieldwork costs and designed sampling scheme using input from ecology and hydrology coursework
    \item Wrote grant abstracts, applications, and reports (awarded \$8,000), and submitted receipts
    \item Conducted data analysis on Icelandic coastal patchy wetlands and evaluated wetland status of protection and restoration both regionally and globally
    \item Assessed and analyzed soil moisture and infiltration characteristics for spatial and statistical differences (ANOVA, t-tests, linear regression)
    \item Analyzed wetland soil characteristics in laboratory (color, grain size, classification, density, pH, etc.)
    \item Coordinated with Icelandic professors to share map of wetland data and findings
    \item Developed maps using ArcGIS, ArcticDEM data, and GPS points for presentations and defense
    \item Researched wetland types for comparison (northern, coastal, volcanic soils, peat) and the impacts of wetland drainage along societal implications on surface hydrology and biology such as migratory birds
    \item Presented results at national conferences, contributed to publications, and shared on social media
    \item Planned logistics for arctic fieldwork (sampling locations, shelter and supplies, RV rental, schedule)
\resumesublistend

\subsectionpositiondate
    {Juneau Icefield Research Program}
    {Research Participant}
    {Jun 2015 - Aug 2015}
\resumesublistbegin
    \item Location: Juneau, AK and Atlin, BC
    \item Advisor: Dr. Allen Pope
    \item Traversed 80+ miles on the Juneau Icefield on an educational expedition living at remote campsites
    \item Conducted glaciological research surveys while learning mountaineering skills and safety, including using a compass, creating an anchor, cross country skiing, and introductory first aid
    \item Attended lectures on glacial geology, geomorphology, hydrology, mass balance, surface energy balance, ecology, indigenous Tlingit oral history, and Alaskan Native relationships with the Juneau Icefield
    \item Gathered information on snow grain size and albedo, using a field spectroradiometer, and analyzed the outcomes of a snow grain size retrieval algorithm
    \item Worked closely on a project team with other students to assess the spectral reflectance and albedo of the Juneau Icefield, focused on testing a snow grain size retrieval algorithm, linking red algae biomass to spectral reflectance, quantifying the radiative forcing of impurities in sun cups, and creating a spectral catalogue of the Taku glacier system
    \item Ensured well-being of others by supporting the group: helped cook for 50+ students/faculty, assisted with set up and break down of campsites and maintained cleanliness, and lead traverse groups or kept pace
\resumesublistend

\section{Current Employment}

\subsectionpositiondate
    {Kaizen Karate}
    {Karate Instructor}
    {Dec 2024 - Present}
\resumesublistbegin
    \item Teach students (40) karate skills and appropriate behavior in class (respect, discipline, motivation)
    \item Check in with schools, answer parent questions, and provide customer service as needed
\resumesublistend

\subsectionpositiondate
    {Hemophilia Foundation of Michigan}
    {Cabin Counselor}
    {Present}
\resumesublistbegin
    \item Attend weekly meetings including planning camp environment and teambuilding activities
\resumesublistend

\section{Employment History}

\subsectionpositiondate
    {Friends of the Detroit River}
    {Stewardship Associate}
    {Nov 2023 - Aug 2024}
\resumesublistbegin
    \item Location: Taylor, MI
    \item Traveled to partner communities to strengthen collaboration with community leaders
    \item Led planning for stewardship program with US Fish and Wildlife Service and Michigan Sea Grant
    \item Oversaw coordination of events at partner institutions and on-site visits, including travel
    \item Translated Watershed Management Plan materials into clear, concise language for general audiences
    \item Communicated with partners and shared environmental monitoring methods and data results
    \item Tracked community outreach and environmental impact progress towards stewardship program goals
    \item Taught master rain gardener course and developed high school curriculum with hands-on exercises
    \item Prepared presentations, lectures, and workshops on invasive species management and conservation
    \item Researched experiential learning strategies, interpretive techniques, and integrated storytelling
    \item Developed rain garden training materials and delivered best practices for rain gardens to professionals
    \item Adapted master rain garden curriculum for high school classrooms
    \item Delivered river cleanup best practices to volunteers and supported Belle Isle Aquarium: tracked 2000+lbs of trash using the CleanSwell app and reported trash types (plastic, biodegradable, styrofoam, misc.)
    \item Proactively engaged with Ecorse Creek Committee stakeholders to encourage cleanup intiatives
    \item Supported operations and communications team coordination to communicate program objectives to interested government institutions, new partners, and to the general public
    \item Developed river conservation training materials and delivered information including: invasive species management, Areas of Concern, and environmental stewardship
    \item Collaborated with Friends of the Detroit River staff and downriver partners to communicate river restoration updates and provide recommendations and guidance for the Ecorse Creek Committee
    \item Strengthened partnerships with federal evaluation experts and non-profit and academic partners
    \item Educated the public on biodiversity indicators (key species, contaminants, flood hazards, habitat loss)
\resumesublistend

\subsectionpositiondate
    {Kids Kicking Cancer (KKC)}
    {Assistant Martial Arts Therapist}
    {Jan 2022 - Present}
\resumesublistbegin
    \item Location: Auburn Hills, MI
    \item Lead elements of class, demonstrate martial arts therapeutic techniques to students, and provide students with a friendly, safe, warm environment to learn
    \item Coordinate hours with other martial arts therapists
    \item Completed online training course in martial arts therapy used in healthcare and check in with supervising therapists for feedback
\resumesublistend

\subsectionpositiondate
    {Native Justice Coalition}
    {Executive Assistant and Development Coordinator}
    {Mar 2023 - Aug 2023}
\resumesublistbegin
    \item Location: Manistee, MI
    \item Supervised team members (4) and addressed general problems and inquiries at events
    \item Advertised and recruited for Board members and communicated with members on behalf of team
    \item Managed grants database, tracked funds, and applied for grants (received \$307,500 in awards)
    \item Organized team meetings, Executive Director's schedule, and planned community events
    \item Administered restorative justice program needs and supported strategic visioning process
    \item Researched Alaska Native community-led programs and Maine Truth and Reconciliation Commission
    \item Determined program priorities and updated program website to align with mission and vision
    \item Developed human resources policies and safety trainings in alignment with mission and vision
    \item Coordinated events with partners and procured lodging, meeting spaces, and supplies
    \item Coordinated and planned restorative justice events for Anishinaabe communities across Michigan
    \item Worked with Executive Director to advertise for and recruit Board members
\resumesublistend

\subsectionpositiondate
    {Juneau Icefield Research Program}
    {Logistics Support Staff}
    {May 2022 - Aug 2022}
\resumesublistbegin
    \item Location: Juneau, AK
    \item Participated in Celebration festival in support and education of Alaska Native community and culture
    \item Mentored and prepared students for the Juneau Icefield expedition and managed faculty needs
    \item Supported concurrent Upward Bound program and identified graduate student teaching strengths
    \item Organized incoming/outgoing travel arrangements and administered reimbursements
    \item Coordinated with Directors and field team to supply research and data collection needs
    \item Facilitated reflection meetings, tracked wellness goals, and lead check ins and debriefs
    \item Served as point of contact and organized participants pre-field requirements
    \item Implemented daily and weekly schedules to manage tasks and prioritize participant needs
    \item Lead students on hikes, managed camp, and organized educational and field resources
    \item Completed training for helicopter safety, Wilderness First Aid, and local hazards
\resumesublistend

\subsectionpositiondate
    {Soil and Land Use Technology, Inc.}
    {Field Technician}
    {Dates Not Provided}
\resumesublistbegin
    \item Managed crew and oversaw collection of soil borings, utilities location, and pavement condition surveys
    \item Gathered soil descriptions, information on infiltration rates and saturated hydraulic conductivity, and tested soils in-situ for volatile organic compounds
    \item Conducted laboratory tests on soils for gravimetric moisture content specific gravity, hydrometer, gradation, and proctor tests
    \item Organized and input soils and pavement condition data for project manager reports
    \item Worked closely in coordination with the Department of General Services lead-in-water testing initiative in D.C. public schools (DCPS) (116): as a liaison between field and office, and assisted in maintaining the Salesforce database for water samples (1,000+), aerators, filters, and testing locations
    \item Worked with teams to implement a water sampling grid around DCPS visiting each quadrant of D.C., oversaw group problem solving, assigned sampling locations, coordinated availabilities
    \item Ensured collection, organization and transportation of water samples  and proper chain of command
    \item Drafted and edited briefings, reports, and memos on quality control and quality analysis
    \item Made data corrections and updates with close attention to detail
    \item Adapted to in-field and office procedural issues
\resumesublistend

\subsectionpositiondate
    {Congress for the New Urbanism}
    {Program Support}
    {Sep 2014 - Jun 2015}
\resumesublistbegin
    \item Location: Chicago, IL
    \item Compiled and analyzed survey data of topic ideas for annual spring congress
    \item Responded to incoming details and plans for speakers throughout the program
    \item Maintained and edited the annual congress website, and oversaw mailing lists of prospective attendees
    \item Managed data from previous congresses and similar organizations for future recommendations
    \item Drafted briefings, presentations and reports, and memos of previous campaign ideas on sustainable urbanism
    \item Campaigned and marketed to gather attendees for the spring congress
    \item Researched biographies of potential board members to organize strong candidates and enable diversity
\resumesublistend

\subsectionpositiondate
    {Resource Center Chicago}
    {Community Intern}
    {Jul 2014 - Nov 2014}
\resumesublistbegin
    \item Built the foundation for urban farms in vacant lots; increased awareness of, and access to, healthy foods in distressed communities on the South Side
    \item Collaborated with other interns to educate local community members of the opportunity for job opportunities through small urban farms
    \item Researched and designed a water conservation system using sprinklers and drip irrigation
    \item Engaged the community (parents, young children) through neighborhood campaigns with flyers and door-to-door canvassing, and assisted with public relations at local farmers markets and art events
\resumesublistend

\subsectionpositiondate
    {Special Events Chicago}
    {Operations Assistant}
    {Jun 2013 - Sep 2013}
\resumesublistbegin
    \item Managed events through direction of services, scheduling, and compliance with visitors
    \item Worked with other team members to assist vendors and other suppliers for a smooth event
    \item Maintained the operations center by checking in visitors and supplying equipment
    \item Facilitated organization and communication by acting as liaison between vendors and management
\resumesublistend

\subsectionpositiondate
    {DePaul University}
    {Education and Development Grant Fellow}
    {Sep 2011 - May 2012}
\resumesublistbegin
    \item Evaluated social media marketing strategies to promote university workshops (career, leadership, financial)
    \item Participated in workshops and courses designed to help students plan career goals and/or graduate studies
    \item Planned quarterly fundraising events to raise awareness of resources for a broad and diverse group of students
\resumesublistend

\section{Additional Employment}
\resumesublistbegin
    \item Walgreens
    \item McDowell \& Associates
    \item Vector Marketing
    \item Outreach
\resumesublistend

\section{Community Leadership and Service}

\subsectionpositiondate
    {Polar Impact}
    {Co-organizer, Social Media Coordinator}
    {Jan 2022 - Present}
\resumesublistbegin
    \item Publish Features project stories to website and co-produce stories by Arctic Indigenous experts
    \item Co-author newsletters and design graphics for social media posts about polar experts
    \item Led all aspects of COP27 social media takeovers: recruitment, advertisement and technical support
    \item Recruit undergraduates, graduates, post-doctoral researchers and experts from other countries to highlight
    \item Share  Features through website and social media platforms including Facebook, Instagram, and X
    \item Schedule meetings taking various time zones into consideration
    \item Co-administer newsletters and design graphics and social media posts
    \item Led COP27 takeovers: coordinated dates, recruitment, advertisement, and technical support
\resumesublistend

\subsectionpositiondate
    {Juneau Icefield Research Program}
    {Inclusion and Diversity Coalition Member}
    {Oct 2020 - Mar 2024}
\resumesublistbegin
    \item Updated Code of Conduct to reflect wellness practices and belonging/inclusion needs in the field
    \item Co-developed action plan to meet inclusion/diversity goals and values
    \item Facilitated and mediated meetings, summarized activity status, results, and new tasks
    \item Developed outreach brochure relating program experience for diverse audiences
    \item Communicated relevant field work insights to issues relating to inclusion and diversity
    \item Developed an updated code of conduct to address accessibility issues
    \item Researched DEI frameworks, wellness practices and adaptable program resources
    \item Updated curriculum materials on climate change and Alaskan Tlingit history
    \item Scheduled and facilitated meetings, summarized actionable items, results, and new tasks
    \item Discussed long term goals and inclusivity issues within the organization
\resumesublistend

\subsectionpositiondate
    {Cranbrook Institute of Science}
    {Program Facilitator}
    {Oct 2020 - Present}
\resumesublistbegin
    \item Location: Bloomfield Hills, MI
    \item Research and develop project-based learning activities to teach various scientific topics at many educational levels for visitors to the Institute (ranging from preschool to elderly ages)
    \item Organize and present interactive materials to engage the community about subject matter relevant to our local habitats and watersheds
    \item Coordinate with employees to set up interactive tables and design materials used
\resumesublistend

\subsectionpositiondate
    {Association of Polar Early Career Scientists}
    {Project Organizer}
    {Oct 2020 - Oct 2021}
\resumesublistbegin
    \item Facilitated and organized interviews of scientists and artists for Antarctica Day 2020
    \item Designed online surveys and collected illustrations and photographs for photo and art competitions
    \item Coordinated Council meetings across international time zones, created project timeline/schedule and implemented activities
\resumesublistend

\subsectionpositiondate
    {Girls on Ice}
    {Application Reviewer}
    {Jan 2019 - Apr 2021}
\resumesublistbegin
    \item Location: Toronto, ON
    \item Reviewed applications to accept participants for Girls on Ice, an expedition-based program to teach young girls about science and art while immersed in nature
    \item Shared and compared quick, thorough decision-making processes with other application reviewers
    \item Maintained confidentiality of young participants and files
    \item Attended volunteer orientation to learn about participation requirements and checking for bias regarding socioeconomic background
\resumesublistend

\subsectionpositiondate
    {DePaul Alumni Sharing Knowledge (ASK)}
    {Mentor}
    {Oct 2017 - Present}
\resumesublistbegin
    \item Remote
    \item Respond to requests for informational interviews on career paths, graduate school
\resumesublistend

\subsectionpositiondate
    {Association of Polar Early Career Scientists (APECS)}
    {Council Member}
    {Oct 2020 - Oct 2021}
\resumesublistbegin
    \item Remote
    \item Coordinated and participated in Council meetings, created project plans for the Arts project group and implemented subgroup activities (workshops, webinars, artists database) to meet goals for the year
    \item Worked with project group team members to design, complete, and track tasks for each subgroup
    \item Prepared for, facilitated, and participated in webinars interviewing Polar scientists in celebration of Antarctica Day, 2020 (December 1st)
    \item Drafted reports, emails, and briefings to connect with experts in Polar science about collaborating on Antarctica Day 2020 webinars
    \item Collected video interview clips and drafted public flyers to share with the general public on different types of Antarctic careers and the scientists' or artists' motivations for working in Antarctica
    \item Drafted online surveys and collected illustrations and photographs for Antarctica Day 2020 photo and art competition
\resumesublistend

\subsectionpositiondate
    {Integra Young Warriors}
    {Assistant Instructor}
    {Sep 2018 - Jan 2020}
\resumesublistbegin
    \item Location: Toronto, ON
    \item Assisted teaching youth, aged 8-11, emotional regulation skills through aikido, yoga, and meditation (mindfulness, reflective listening, and breathing skills to focus in on the mind-body relationship)
    \item Led martial arts exercises and provided space for 1-on-1 meets between students and instructors
    \item Created an empathetic, patient, fun space; taught the importance of focusing, finding balance, and helped reinforce learned skills for more advanced students
    \item Participated in orientation about Learning Disabilities/Mental Health and components of psychoeducation to understand case-specific challenges, including Walk-a-Mile exercises
\resumesublistend

\subsectionpositiondate
    {Hunt Valley Symphony Orchestra (HVSO) | Baltimore Philharmonia Orchestra (BPO)}
    {Violoncello Performer}
    {Apr 2016 - Jun 2017}
\resumesublistbegin
    \item Location: Baltimore, MD
    \item Performed the violoncello with community-based, non-profit, volunteer orchestras
    \item Provided a musical outlet through performances, promoting music education in Hunt Valley and accessibility to public concerts in the Baltimore area at no cost with the BPO
    \item Built community relationships with school music directors and composers through collaboration with other musician members and advertised local opportunities to hear performances
\resumesublistend

\subsectionpositiondate
    {Back On My Feet}
    {Volunteer}
    {Nov 2016 - Jun 2017}
\resumesublistbegin
    \item Location: Baltimore, MD
    \item Attended orientation to learn appropriate behavior and  interactions with those experiencing homelessness within the Back On My Feet community
    \item Participated in weekly runs with other volunteers and residents of transitional shelters to help occupants remain accountable and continue to achieve progressive goals and stability
\resumesublistend

\section{Additional Volunteer Activities}
\resumesublistbegin
    \item Muslim Women Center
    \item Volunteering at the Children's Center in Detroit
    \item Volunteering at the Capuchin soup kitchen
\resumesublistend

\section{Awards}
\resumesublistbegin
    \item \textbf{2023:} Michigan Voices Capacity Grant -- Michigan Truth \& Reconciliation Commission (\$7,500)
    \item \textbf{2023:} WK Kellogg Foundation -- Native Justice Coalition, General Operating Support (\$300,000)
    \item \textbf{2018:} York University -- Fieldwork Costs Fund (\$2,600)
    \item \textbf{2018:} American Scandinavian Foundation Fellow -- Scandinavian Seminar (\$5,000)
    \item \textbf{2018:} Society of Wetland Scientists -- Student Research Grant (\$1,000)
    \item \textbf{2017-2018:} York University -- Graduate Scholarship (\$6,000)
    \item \textbf{2011-2015:} St. Vincent DePaul -- Academic Scholarship (\$24,000)
    \item \textbf{2013-2014:} Sociology -- Endowed Student Scholarship (\$3,000)
    \item \textbf{2011-2012:} DePaul -- Education and Development Grant for Employability Award (\$3,000)
\resumesublistend

\section{Publications and Presentations}

\subsectiondate
    {2024}
    {3}
\resumesublistbegin
    \item Dryák-Vallies, M., Gill, P., Greenly, T., Haycock-Chavez, N., Nelson, M., and Perera, E. (2024), Restoring Stewardship // Unangam Maqax\^{}singis suug\^{}utakus, article, retrieved from: \href{https://www.polarimpactnetwork.org/lf-restoring-stewardship}{polarimpactnetwork.org/lf-restoring-stewardship}
    \item Perera, E., Leadership Adaptation through an Open Mind, talk presented at Oregon State University, Corvallis, OR, 24 Oct
    \item Perera, E., Walsh, S., Master Rain Gardener High School Series, presentation series given during Friends of the Detroit River sponsored courses at Grosse Ile High School, Grosse Ile, MI, Feb--Jun
    \item Perera, E., Soils; Infiltration and Percolation, Southeast Michigan Master Rain Garden Class, virtual presentations given during Master Rain Garden webinar series, Detroit, MI, Jan--Feb
\resumesublistend

\subsectiondate
    {2021}
    {1}
\resumesublistbegin
    \item Miller, R., Fortner, S., Lakeram, S., Leidman, S., Peek, M., and Perera, E. (2021), Unlearning Racism in Geosciences: A Case Study from the Juneau Icefield Research Program, virtual poster presented at 2021 AGU Fall Meeting, retrieved from: \href{https://agu2021fallmeeting-agu.ipostersessions.com/Default.aspx?s=0C-7C-95-47-65-84-5B-A9-3A-75-41-46-54-B9-C8-18}{agu2021fallmeeting-agu.ipostersessions.com}
\resumesublistend

\subsectiondate
    {2019}
    {2}
\resumesublistbegin
    \item Young, K.L., Scheffel, H.-A. and Perera, E.A. (2019), Impact of dust on the local hydrology of Arctic Landscapes. Paper presented at 22nd International Northern Research Basins Symposium and Workshop, Yellowknife, NWT, 18-23 Aug
    \item Perera, E.A. and Young, K.L. (2019), Initial surface hydrology characteristics of Icelandic, drained, patchy wetlands. Wetland Science \& Practice, Apr: 83-90
\resumesublistend

\subsectiondate
    {2018}
    {2}
\resumesublistbegin
    \item Perera, E.A., and Young, K.L. (2018), Variation in Infiltration and Soil Moisture of Icelandic Drained, Patchy Wetlands, Poster (H43E-2441) presented at 2018 AGU Fall Meeting, Washington, DC, 10-14 Dec
    \item Perera, E.A. and Young, K.L. (2018), Infiltration of Icelandic Drained, Patchy Wetlands, Poster (ABS717) presented at 2018 IGU-CAG-NCGE Conference, Québec City, QC, 6-10 Aug
\resumesublistend

\subsectiondate
    {2015}
    {3}
\resumesublistbegin
    \item Hughes-Allen, L., Peter, A., Perera, E., Pope, A., and Popyack, K. (2015), Reflectance Spectra of the Juneau Icefield, Poster (C13A-0794) presented at 2015 AGU Fall Meeting, San Francisco, CA, 14-18 Dec
    \item Hughes-Allen, L., Peter, A., Perera, E., Pope, A., and Popyack, K., Reflectance Spectra of the Juneau Icefield, talk given at Juneau Icefield Research Program 2015 Research Presentations, Juneau, AK
    \item Hughes-Allen, L., Peter, A., Perera, E., Pope, A., and Popyack, K., Reflectance Spectra of the Juneau Icefield, talk given at Juneau Icefield Research Program 2015 Research Presentations, Atlin, BC
\resumesublistend

\section{Languages and Certifications}
\resumesublistbegin
    \item April 2025 -- Adult Mental Health First Aid
    \item May 2024 -- Adult and Pediatric First Aid/CPR/AED
    \item ESP -- Conversational Spanish (A2/B1)
\resumesublistend

\section{Instrumentation Used}
Double ring infiltrometer, HOBO data loggers and 10HS moisture sensor, OceanOptics USB4000 spectrometer, FLIR thermal imaging camera, infrared gas analyzer (IRGA), Aardvark permeameter, pneumatic consolidometer, photoionization detector (PID), ASD field spectroradiometer, various in-situ methods. Competencies in HTML, SPSS, MATLAB, ArcGIS and MicroStation.

\section{Class Projects}

\subsectiondate
    {Soil quality assessment}
    {}
\resumesublistbegin
    \item Collected soil profile descriptions on an urban farm and a suburban neighborhood for comparison
    \item Assessed soil infiltration, soil compaction, and soil temperature
    \item Conducted laboratory tests to assess Carbon, Nitrogen, Phosphorus, bulk density, and soil texture
    \item Wrote a report on land history, research methodology, and detailed findings and implications
\resumesublistend

\subsectiondate
    {Political economy analysis of oil drilling in Alaska}
    {}
\resumesublistbegin
    \item Gathered information on the history of oil drilling pros and cons and the industry relationship to Alaskans and the Alaskan economy
    \item Assessed whether oil drilling would be viable in the Arctic National Wildlife Refuge using a Cost-Benefit Analysis model
    \item Investigated background of oil drilling in Ecuador for comparison to the Arctic National Wildlife Refuge
    \item Examined federal and state law history to determine specific benefits of oil drilling to Alaska and the lower 48 States
\resumesublistend

\subsectiondate
    {Senior capstone: Veteran homelessness}
    {}
\resumesublistbegin
    \item Reviewed 2015 policy implications on veterans experiencing homelessness, including veteran benefits or the lack thereof within the Department of Housing and Urban Development and the Department of Veterans Affairs
    \item Updated policies from 2012 paper on veterans experiencing homelessness in the Chicagoland area, starting from a macro to a micro-level scale, including state run resources provided compared to city level resources provided
    \item Investigated the potential of holistic solutions and outcomes for those veterans, such as Housing First
\resumesublistend

\end{document}
